\documentclass[10pt,letterpaper]{article}

% TeXShop / TeX Live 2019-friendly preamble
\usepackage[T1]{fontenc}
\usepackage[utf8]{inputenc}

\usepackage[margin=0.5in]{geometry}
\usepackage{multicol}
\usepackage{enumitem}
\usepackage{xcolor}
\usepackage{hyperref}
\usepackage{fancyhdr}

\setlength{\parindent}{0pt}
\setlength{\parskip}{1.5pt}
\setlist[itemize]{leftmargin=*, itemsep=1pt, topsep=1.5pt, parsep=0pt, partopsep=0pt}

\definecolor{Banner}{RGB}{20,20,20}
\definecolor{LightGray}{gray}{0.92}
\definecolor{RuleGray}{gray}{0.78}

% Footer (does not affect column flow)
\pagestyle{fancy}
\fancyhf{}
\renewcommand{\footrulewidth}{0.35pt}
\fancyfoot[C]{\footnotesize ACO v1.2 \hspace{0.6em}\textbullet\hspace{0.6em} MedSecurance / O-ETB \hspace{0.6em}\textbullet\hspace{0.6em} Reference Card}
\renewcommand{\headrulewidth}{0pt}

\newcommand{\CardTitle}[1]{%
  \begin{center}
    {\Large\bfseries #1}\\[-2pt]
    \rule{0.92\linewidth}{0.6pt}
  \end{center}
}

\newcommand{\BannerBox}[1]{%
  \noindent\colorbox{Banner}{%
    \parbox{\linewidth}{\color{white}\bfseries\small\hspace{6pt}#1\hspace{6pt}}%
  }\\[10pt]
}

% Divider BEFORE header
\newcommand{\Sec}[1]{%
  \vspace{5pt}%
  \noindent{\color{RuleGray}\rule{\linewidth}{0.35pt}}\par
  \vspace{2pt}%
  {\bfseries #1}\par
  \vspace{2pt}%
}

% Code-style example box with preserved indentation
\newcommand{\Example}[1]{%
  \vspace{1pt}%
  \noindent\colorbox{LightGray}{%
    \parbox{\linewidth}{\ttfamily\footnotesize\raggedright
#1
    }%
  }%
  \vspace{2pt}%
}

\newcommand{\DoDont}[2]{%
  \vspace{1pt}%
  \noindent\colorbox{LightGray}{%
    \parbox{\linewidth}{\small\textbf{DO:}~#1\hfill \textbf{DON'T:}~#2}%
  }%
  \vspace{2pt}%
}

% Reduce overfull boxes
\sloppy

\begin{document}

% ===================== PAGE 1 =====================
\CardTitle{Assurance Case Outline (ACO) v1.2 Quick Reference Card}
\BannerBox{AUTHORS}

\begin{multicols*}{2}
\raggedcolumns

\Sec{Purpose}
\begin{itemize}
  \item Lightweight, indentation-based outline language for assurance cases.
  \item Primary authoring interface to O-ETB; one source drives APL + downstream artefacts.
  \item Designed to be readable, unambiguous, and tooling-friendly.
\end{itemize}
\Example{Case: Operational Plane --- Policy scope\\
Goal G0 operationalPlane:\\
\hspace*{1.5em}The operational plane guarantees the \string\{local policy\string\} is met.}

\Sec{Node Types (Core Set)}
\begin{itemize}
  \item \textbf{Goal} -- claim that must hold.
  \item \textbf{Strategy} -- explains how children jointly support a claim.
  \item \textbf{Context} -- facts, definitions, scope/environment.
  \item \textbf{Assumption} -- condition taken as given.
  \item \textbf{Justification} -- rationale for acceptability.
  \item \textbf{Evidence} -- artefact that discharges a claim (leaf).
  \item \textbf{Module} -- encapsulated/reusable argument fragment (goal-like).
\end{itemize}

\Sec{Header Syntax (Mandatory)}
\Example{<Type> [ID] <shortLabel>:}
\begin{itemize}
  \item Header lines must end with a colon \texttt{:}.
  \item \texttt{Type} is one of: Goal, Strategy, Context, Assumption, Justification, Evidence, Module.
  \item \texttt{ID} is optional (single token; e.g., \texttt{G0}, \texttt{C2}, \texttt{M1}).
  \item \texttt{shortLabel} is either:
  \begin{itemize}
    \item unquoted single token (\texttt{componentBehaviour}), or
    \item quoted short text (\texttt{"Component Behaviour"}).
  \end{itemize}
\end{itemize}
\Example{Goal G1 componentBehaviour:\\
Context "Plane Definition":\\
\hspace*{1.5em}Behaviour of the \string\{component\string\} satisfies the local policy.}

\Sec{Body Text}
\begin{itemize}
  \item One or more natural-language lines indented beneath the header.
  \item May span multiple lines; blank lines allowed for readability.
\end{itemize}
\Example{Goal G1 componentBehaviour:\\
\hspace*{1.5em}Behaviour of the \string\{component\string\} satisfies the local policy.\\
\hspace*{1.5em}This may be justified by component test evidence.}

\columnbreak

\Sec{Indentation Rules (Structure)}
\DoDont{Increase indentation by one level at a time.}{Jump levels (e.g., level 0 -> level 3).}
\begin{itemize}
  \item Recommended: 2 spaces per level; tabs discouraged.
  \item A node at level \texttt{k} is child of nearest preceding node at level \texttt{k-1}.
  \item Indentation defines structure and default semantics.
\end{itemize}

\Sec{Default Semantics (Automatic)}
\begin{itemize}
  \item Child \textbf{Context} \\
        \texttt{=> in\_context\_of(parent, child)}.
  \item Child \textbf{Goal / Strategy / Module / Evidence / Assumption / Justification} \\
        \texttt{=> supported\_by(parent, child)}.
\end{itemize}

\Sec{Evidence Rules}
\DoDont{Attach Context/Assumption/Justification related to evidence to the Evidence's parent.}{Nest any node under an Evidence node.}
\begin{itemize}
  \item Evidence nodes are leaf-level dischargers.
  \item Provisional evidence is permitted; tools distinguish provisional vs.\ accepted.
\end{itemize}
\Example{Goal G2 cryptoCorrectness:\\
\hspace*{1.5em}Cryptography is correctly implemented.\\
\hspace*{1.5em}Evidence E0 Planned "Test report":\\
\hspace*{3.0em}Planned: vendor conformance test results (to be produced).}

\Sec{Modules (Reusable Fragments)}
\begin{itemize}
  \item Written using standard header syntax: \texttt{Module [ID] label:}
  \item Behave like specialised Goals: may contain internal sub-goals/strategies/evidence.
  \item Use for reuse, decomposition by component/subsystem, and scale management.
\end{itemize}
\Example{Module M1 component:\\
\hspace*{1.5em}Behaviour of the \string\{component\string\} meets its local policy.}

\Sec{Optional Syntax}
\begin{itemize}
  \item \textbf{Case header (metadata):} \texttt{Case: <Title> - <Scope>} (optional; not a node).
  \item \textbf{Explicit relations (optional):} natural-language relation sentences (see page 2).
  \item \textbf{Placeholders:} use braces \texttt{\string\{...\string\}} and define them nearby or in a Context.
\end{itemize}

\end{multicols*}

\clearpage

% ===================== PAGE 2 =====================
\CardTitle{Assurance Case Outline (ACO) v1.2 Quick Reference Card}
\BannerBox{REVIEWERS / ASSESSORS}

\noindent\colorbox{LightGray}{%
\parbox{\linewidth}{\small
\textbf{Fast scan:}\hspace{0.6em}[ ] Top goals supported \hspace{0.6em}\hspace{0.6em}[ ] Evidence closes claims \hspace{0.6em}\hspace{0.6em}[ ] Assumptions justified \hspace{0.6em}\hspace{0.6em}[ ] No structural errors
}}\\[6pt]

\begin{multicols*}{2}
\raggedcolumns

\Sec{What Tools Derive Automatically}
\begin{itemize}
  \item Node Model (nodes + relations) from headers, bodies, indentation, and explicit relation sentences.
  \item Canonical hierarchical IDs computed from indentation (for diagnostics/canonicalisation).
  \item ASCII tree views at multiple verbosity levels; plus summary statistics.
\end{itemize}

\Sec{Errors (Must Fix)}
\begin{itemize}
  \item Duplicate node IDs.
  \item Dangling references in explicit relation sentences.
  \item Malformed headers or unknown node types.
  \item Indentation jumps (skipping levels).
  \item Evidence nodes with children.
\end{itemize}

\Sec{Warnings / Quality Signals}
\begin{itemize}
  \item \textbf{Undeveloped} goals/modules (insufficient supporting structure).
  \item Provisional evidence occurrences.
  \item Undefined or inconsistent placeholders.
  \item Ambiguous wording in Goal/Module bodies.
\end{itemize}

\Sec{v1.2 State Distinctions}
\begin{itemize}
  \item \textbf{Undeveloped:} concrete node whose argumentation is incomplete.
  \item \textbf{Uninstantiated:} refers to a future or unresolved entity.
  \item These states are distinct and may co-occur.
\end{itemize}
\Example{Undeveloped: Goal has no sub-goals/evidence.\\
Uninstantiated: Module reference exists but target module not imported/resolved.}

\Sec{Explicit Natural-Language Relations (Optional)}
\begin{itemize}
  \item \textbf{Supported-by:} \texttt{G0 is supported by G1, G2 and M0.}
  \item \textbf{In-context-of:} \texttt{C0 provides context for G0.}
  \item Explicit relations augment (not replace) indentation-derived relations.
\end{itemize}
\Example{G0 is supported by G1, G2 and M0.\\
C0 provides context for G0.}

\columnbreak

\Sec{ASCII Tree Review}
\begin{itemize}
  \item \textbf{Skeleton:} headers only.
  \item \textbf{Structure:} brief body snippet.
  \item \textbf{Full:} complete body text.
\end{itemize}

\Sec{Canonicalisation}
\begin{itemize}
  \item Rewrites headers and relation sentences to use canonical hierarchical IDs.
  \item Recommended once structure settles.
\end{itemize}

\Sec{Reviewer Heuristics}
\begin{itemize}
  \item Every top-level Goal has explicit support?
  \item Claims closed by accepted evidence?
  \item Assumptions visible and justified?
  \item Excessive provisional evidence?
\end{itemize}

\Sec{Pipeline Context}
\Example{ACO -> Node Model -> APL -> Instantiate -> Reports/Graphs}

\end{multicols*}

\end{document}